\documentclass[12pt]{article}

\usepackage{graphicx}
\usepackage{amsmath}
\usepackage{geometry}
\usepackage{xcolor}
\geometry{margin=1in}

\begin{document}

%================ PAGE 1 ===================

\begin{center}

\includegraphics[width=2.5cm]{iiitb_comet_logo.jpg}
\end{center}

\vspace{0.2cm}

\begin{center}
\textbf{\large Srinidhi A} \\[3pt]
\textbf{COMETFWC053}
\end{center}

\vspace{0.4cm}

\begin{center}
{\Large \textbf{GATE Question No. 41}}
\end{center}

\vspace{0.3cm}
\section*{\color{blue}Question}
\noindent
For any set of inputs $A$ and $B$, the following circuits give the same output $Q$, except one. Identify the circuit which gives a different output. (GATE PH 2010)

\vspace{0.4cm}

\begin{center}
\includegraphics[width=9.5cm]{question.jpg}
\end{center}

\newpage

%================ PAGE 2 ===================

\section*{\color{blue}Question Analysis}

From logical simplification of the circuits, the Boolean expression obtained is:

\[
Q = \overline{A} + B
\]

Each circuit is verified against this expression to determine the incorrect one.

%================ TRUTH TABLE ===================

\section*{\color{blue}The Truth Table}

\begin{center}
\begin{tabular}{|c|c|c|}
\hline
A & B & $Q = \overline{A} + B$ \\
\hline
0 & 0 & 1 \\
0 & 1 & 1 \\
1 & 0 & 0 \\
1 & 1 & 1 \\
\hline
\end{tabular}
\end{center}

%================ HARDWARE ===================

\section*{\color{blue}Hardware Implementation}

The Boolean expression is implemented using the Arduino UNO board. Inputs $A$ and $B$ are provided through digital pins and the output $Q$ is verified using an LED / 7-segment display.

%================ COMPONENTS ===================

\section*{\color{blue}Required Components \& Pin Connections}

\begin{center}
\begin{tabular}{|c|l|}
\hline
\textbf{Component} & \textbf{Connection} \\
\hline
Input A & Digital Pin 2 \\
Input B & Digital Pin 3 \\
Output Q & Digital Pin 7 \\
GND & GND \\
VCC & 5V \\
\hline
\end{tabular}
\end{center}

\vspace{0.3cm}

\textbf{Components Used:}

\begin{itemize}
\item Arduino UNO
\item Breadboard
\item LED / 7 Segment Display
\item Resistors
\item Jumper Wires
\item USB Cable
\end{itemize}

%================ LOGIC DESCRIPTION ===================

\section*{\color{blue}Logic Description}

\begin{itemize}
\item Initialize inputs $A$ and $B$.
\item Implement the expression $Q = \overline{A} + B$.
\item Apply all input combinations as per truth table.
\item Verify hardware output with theoretical values.
\end{itemize}

%================ CODE STEPS ===================

\section*{\color{blue}Code Uploading Steps}

\begin{enumerate}
\item Create a Platform IO project.
\item Write the code in \texttt{main.cpp} inside the \texttt{src} folder.
\item Run the PIO project using the command \texttt{pio run}. This compiles the code and generates the .hex file.
\item Copy the .hex file to the Arduino Droid folder.
\item Connect the Arduino UNO to the mobile using an OTG cable.
\item Upload the hex file using the "Upload Precompiled" option.
\item Observe the output and verify the expression.
\end{enumerate}

%================ EXPERIMENTAL TABLE ===================

\section*{\color{blue}Experimental Truth Table}

\begin{center}
\begin{tabular}{|c|c|c|}
\hline
A & B & Observed Q \\
\hline
0 & 0 & 1 \\
0 & 1 & 1 \\
1 & 0 & 0 \\
1 & 1 & 1 \\
\hline
\end{tabular}
\end{center}

\vspace{0.5cm}

\begin{center}
\includegraphics[width=10cm]{hardware.jpg}
\end{center}

%================ CONCLUSION ===================

\section*{\color{blue}Conclusion}

The Boolean expression $Q = \overline{A} + B$ was successfully implemented using the Arduino UNO board.

The hardware output matches the theoretical truth table. Hence, the experiment has been verified successfully.

\end{document}
